% Generated draft: requires mathematical and source review.
\documentclass[11pt]{article}
\usepackage[margin=25mm]{geometry}
\usepackage{fontspec}
\setmainfont{Latin Modern Roman}
\usepackage{amsmath,amssymb,mathtools,unicode-math}
\setmathfont{Latin Modern Math}
\usepackage{enumitem,needspace,xurl,hyperref}
\setlength{\parindent}{0pt}
\setlength{\parskip}{5pt}
\setlength{\emergencystretch}{3em}
\newcommand{\N}{\mathbb{N}}
\newcommand{\Z}{\mathbb{Z}}
\newcommand{\Q}{\mathbb{Q}}
\newcommand{\R}{\mathbb{R}}
\newcommand{\C}{\mathbb{C}}
\newcommand{\F}{\mathbb{F}}
\newcommand{\A}{\mathbb{A}}
\newcommand{\PP}{\mathbb{P}}
\newcommand{\E}{\mathbb{E}}
\newcommand{\eps}{\varepsilon}
\newcommand{\dd}{\,\mathrm{d}}
\begin{document}
\section*{UnsolvedMath ID: 30004264}
\textbf{Source problem code:} catalogue entry
\subsection{Definitions and mathematical statement}
For a quiver \(Q\), a representation \(M\) over a field \(k\), and a dimension vector \(e=(e_v)_{v\in Q_0}\), the quiver Grassmannian \(\operatorname{Gr}_e(M)\) parametrizes subrepresentations \(N\subseteq M\) satisfying \(\dim_k N_v=e_v\) at every vertex \(v\). The group \(\operatorname{Aut}_Q(M)\) of invertible morphisms of representations acts by \(N\mapsto gN\). Under an invariant-function interpretation, invariance means \(f(gN)=f(N)\) for every \(g\in\operatorname{Aut}_Q(M)\) and every point \(N\) in the chosen domain of the function.

The supplied record specifies the space
\[
 A=\bigoplus_{[P],[I],e}
 \mathbb{Q}^{\operatorname{Aut}_Q(P\oplus I)}
 [\operatorname{Gr}_e(P\oplus I)].
\]
Here \([P]\) and \([I]\) denote isomorphism classes of representations in the intended indexing families. Although the notation suggests projective representations \(P\) and injective representations \(I\), those families and the meaning of ``principal'' are not defined in the supplied record. The notation for each invariant-function summand is also retained without choosing between all functions, constructible functions, or another function space.

\textbf{Requested construction.} Define a Hall-algebra-type multiplication on \(A\) that makes it a graded associative algebra and induces an isomorphism
\[
 \Phi:U_q(\mathfrak b^+)\otimes U_q(\mathfrak b^+)
 \longrightarrow A.
\]
Under a \(\mathbb Q\)-vector-space interpretation of the summands, the requested multiplication would be a bilinear map \(\mu:A\otimes_{\mathbb Q}A\to A\) satisfying
\[
 \mu(\mu(a,b),c)=\mu(a,\mu(b,c))
 \qquad(a,b,c\in A),
\]
and, once a grading \(A=\bigoplus_\gamma A_\gamma\) and its index operation are specified,
\[
 \mu(A_\gamma,A_\delta)\subseteq A_{\gamma+\delta}.
\]
The symbol \(U_q(\mathfrak b^+)\) is the intended quantum enveloping algebra of a positive Borel subalgebra; neither that subalgebra nor its relation to \(Q\) is specified. The class of quivers, field, indexing families, function theory, grading, parameter \(q\), coefficient ring and tensor-product base, Hall-type compatibility conditions, construction of \(\Phi\), and required category of isomorphism remain unresolved. Thus these formulas record the construction problem but do not supply a fully determined existence statement or a universal quantifier over quivers.
\subsection{Short English statement}
The problem asks for a Hall-algebra-type product on invariant functions on principal quiver Grassmannians, giving an associative graded algebra identified with two tensor factors of \(U_q(\mathfrak b^+)\). The supplied record leaves essential hypotheses and the meaning of this identification unspecified; these require mathematical review.
\subsection{Sources}
\begin{itemize}
\item Local UnsolvedMath registry, \texttt{data/unsolvedmath/problems.json}, permanent ID \(30004264\), title ``Hall Multiplication on Principal Quiver Grassmannians,'' source problem code \texttt{catalogue entry}, in the Oberwolfach Reports Open Problems collection. The record lists proposed year 2019 and names no individual proposer.
\item Supplied source citation: \emph{Mini-Workshop: Degeneration Techniques in Representation Theory (2020)}, Oberwolfach Reports, DOI \texttt{10.4171/owr/2019/46}, pp.\ 2869--2909. \url{https://doi.org/10.4171/owr/2019/46}. The supplied background describes Oberwolfach Reports as documenting mathematical workshops and their open-problem discussions.
\item Supplied literature source: ``OWR 2019 46,'' \url{https://ems.press/content/serial-article-files/46825}.
\item Literature assessment dated exactly \texttt{2026-08-22}: the supplied curation labels the problem open and reports that the proposed Hall-algebra-type multiplication is posed as a problem and that no general construction with the stated properties was found. This is supplied dated curation derived from the OWR record, not an independently verified literature assessment; time-sensitive status and cited claims require independent verification.
\end{itemize}
\end{document}
